\section{Limitations and what is next}
\label{sec:discussion}

Our tracking pipeline refines tile positions from a single session using median single-tile poses rather than full bundle adjustment \citep{triggs2000bundle}; joint optimization over all multi-tile observations should further reduce single-tile jitter below the current \SI{4}{}--\SI{9}{mm}.
Because 6-DoF tracking requires simultaneous visibility of the glove and at least one reference tile, label coverage depends on tile placement; while we establish empirical placement rules (\S\ref{sec:system}), we did not evaluate alternate layouts.
On the robot, the camera positioner lags gripper motions by \SI{0.4}{s} and can reach joint limits during wide sweeping movements. A faster positioner arm or lighter camera payload would reduce this latency and expand the tracking workspace. The viewpoint model itself received little effort: a linear fit on one session, with no tuning against rollouts, and the rollouts still showed the jaws frequently occluded by the robot's own wrist-camera bracket, which stayed on the gripper so that the wrist-view arms could be run on the same days. Removing that bracket, as the design intends, and fitting the model on more sessions are the cheapest next gains for the head-camera arms.

An appearance gap also persists between human demonstrations and robot deployment: the head camera captures the collector's forearm and glove fingers, preventing policies trained solely on glove demonstrations from transferring zero-shot (robot validation loss 0.52 vs.\ 0.26 co-trained), as observed in the fixed-camera version of the device.
Several experimental controls also warrant future evaluation.
First, the fixed-camera baseline was mounted low because the positioner occupied the higher vantage point; a higher fixed camera would reduce occlusion, though it still cannot follow the hand.
Second, teleoperators guided the robot while viewing the head-camera preview, which may favor egocentric observations; we did not collect a control set teleoperated under a wrist-camera preview. Third, each arm is a single seed evaluated once per placement, so the budget pair and the fixed-camera pair sit at $p = 0.06$ and the remaining pairs above it: the ranking is consistent across every paired comparison, but the individual gaps are not yet significant at 25 placements.
Third, although our low-cost design enables parallel collection across multiple operators, all four demonstration sessions were recorded by a single collector with one glove.

Collecting demonstrations in a workspace away from the robot needs no robot-side calibration: the tile-to-base link, fitted here once with the robot holding the ball, can instead be left free and recovered per session, as UMI-style pipelines do with trajectories expressed relative to their start. What the collector must do is lay enough tiles that at least one, preferably two, is in view throughout the task, and tape them down so they do not move within a session.
Hardware designs \citep{yamumi2026} and calibration, recording, and tracking tools \citep{forwardcamumi2026} are open under the MIT license; robot-side hand-eye calibration and dataset-construction utilities are not yet released.
We evaluated a single arm to keep the comparison clean, but nothing in the approach is single-handed: a second glove with its own ball and tag ids is read by the same head camera, which in bimanual work already frames both hands, and on the robot the one positioner serves both arms. Key next steps include that bimanual version, multi-tile workspace mapping, parallel collection across diverse in-the-wild environments, and a second manipulation task.
